\section{Discussion and Conclusion}

scProto addresses two core limitations of existing approaches by combining a community-preserving objective with reconstruction-based cross-batch alignment: affinity-based methods preserve local structure but remain within-batch, limiting coverage of rare populations; reconstruction-based models align across batches but lose community structure, absorbing rare cell states into denser regions of the latent space. Two limitations remain: the number of prototypes $K$ must be specified in advance and suboptimal choices can affect performance, and downstream quality depends on the affinity graph---noisy or poorly constructed affinities propagate misleading structure into the learned prototypes.

The framework extends naturally beyond transcriptomics: spatial affinities enable scProto to capture niche-correlated transcriptional programs while the reconstruction objective preserves cell-type identity. Future directions include more principled spatial affinities based on ligand--receptor interactions or local cell-type composition, and affinity-guided prototypes as compact, interpretable units within single-cell foundation models~\citep{szalata2024transformers,richter2025delineating,bahrami2025scconcept}---aggregating pretrained embeddings while preserving local structure and complementing global alignment in multimodal settings~\citep{richter2026beyond,bahrami2026modality}.
